\documentclass[11pt,a4paper]{article}

\usepackage{amsmath,amssymb,amsthm}
\usepackage{mathtools}
\usepackage{geometry}
\usepackage{enumitem}
\usepackage{hyperref}
\usepackage{cleveref}

\geometry{margin=1in}

\newtheorem{theorem}{Theorem}[section]
\newtheorem{proposition}[theorem]{Proposition}
\newtheorem{lemma}[theorem]{Lemma}
\newtheorem{corollary}[theorem]{Corollary}
\theoremstyle{definition}
\newtheorem{definition}[theorem]{Definition}
\newtheorem{example}[theorem]{Example}
\newtheorem{remark}[theorem]{Remark}
\newtheorem{assumption}[theorem]{Assumption}

\DeclareMathOperator{\Res}{Res}
\DeclareMathOperator{\Ind}{Ind}
\DeclareMathOperator{\rank}{rank}
\DeclareMathOperator{\tr}{tr}
\DeclareMathOperator{\diag}{diag}
\DeclareMathOperator{\spec}{spec}
\DeclareMathOperator{\wind}{wind}
\newcommand{\C}{\mathbb{C}}
\newcommand{\R}{\mathbb{R}}
\newcommand{\Z}{\mathbb{Z}}
\newcommand{\N}{\mathbb{N}}
\newcommand{\E}{\mathbb{E}}
\newcommand{\cL}{\mathcal{L}}
\newcommand{\cH}{\mathcal{H}}
\newcommand{\cD}{\mathcal{D}}
\newcommand{\cW}{\mathcal{W}}
\newcommand{\cN}{\mathcal{N}}
\newcommand{\eps}{\varepsilon}

\title{Meromorphic Structure of Saddle-to-Saddle Learning Dynamics\\in Recurrent Networks}
\author{}
\date{}

\begin{document}
\maketitle

\begin{abstract}
We develop a rigorous framework for understanding the learning dynamics of
recurrent neural networks in terms of the pole-zero configuration of their
transfer functions, viewed as meromorphic functions of a complex variable.
Working with a general time-domain loss under arbitrary input distributions,
we prove that pole-zero cancellations correspond precisely to saddle points
of the population loss---a consequence of the multiplicative
$r_k = \alpha_k\beta_k$ structure of residues, independent of the choice
of loss or input statistics. The transitions between these saddles are
topologically quantised by a winding number invariant derived from the
argument principle, and the sequential activation of dynamical modes during
learning is governed by the \emph{effective spectral energy}---the product
of the target's spectral content and the input's spectral density at each
mode's frequency. This extends the saddle-to-saddle framework of
Saxe et al.\ from deep linear feedforward networks to the recurrent
setting, replacing singular values with the pole-zero geometry of the
complex plane. We derive exact dwell times at each saddle, characterise
the geometry of pole-zero separation in $\C$, and identify mechanisms
for catastrophic interference as non-partner pole-zero collisions.
\end{abstract}

\tableofcontents

%======================================================================
\section{Setup and Notation}\label{sec:setup}
%======================================================================

\begin{definition}[Linear Recurrent System]\label{def:lrs}
A \emph{parametrised linear recurrent system} is a triple
$(W(\tau), B(\tau), C(\tau))$ depending smoothly on a learning
parameter $\tau \in [0,\infty)$, with $W(\tau) \in \R^{n \times n}$,
$B(\tau) \in \R^{n \times d_{\mathrm{in}}}$,
$C(\tau) \in \R^{d_{\mathrm{out}} \times n}$, defining the recurrence
\begin{equation}\label{eq:recurrence}
  h_t = W(\tau)\, h_{t-1} + B(\tau)\, x_t, \qquad
  y_t = C(\tau)\, h_t,
\end{equation}
where $h_t \in \R^n$ is the hidden state,
$x_t \in \R^{d_{\mathrm{in}}}$ is the input, and
$y_t \in \R^{d_{\mathrm{out}}}$ is the output.
The learning timescale $\tau$ is assumed slow relative to the
recurrence index $t$.
\end{definition}

\begin{definition}[Transfer Function]\label{def:transfer}
The $z$-domain transfer function of the system
\eqref{eq:recurrence} is the matrix-valued meromorphic function
\begin{equation}\label{eq:transfer}
  H(z;\tau) = C(\tau)\bigl(zI - W(\tau)\bigr)^{-1}B(\tau).
\end{equation}
\end{definition}

\noindent
Throughout this paper, we work primarily in the SISO
(single-input single-output) setting $d_{\mathrm{in}} = d_{\mathrm{out}} = 1$
to make the pole-zero structure maximally transparent. The MIMO
generalisation is discussed in \Cref{sec:mimo}.

\begin{definition}[Poles, Zeros, and Residues]\label{def:poles_zeros}
In the SISO case, write $B(\tau) = b(\tau) \in \R^n$,
$C(\tau) = c(\tau)^{\top} \in \R^{1 \times n}$. Let
$\{p_k(\tau)\}_{k=1}^n$ be the eigenvalues of $W(\tau)$
(counted with algebraic multiplicity), with corresponding right
eigenvectors $v_k(\tau)$ and left eigenvectors $u_k(\tau)$ normalised
so that $u_k^* v_k = 1$. Then:
\begin{enumerate}[label=(\roman*)]
  \item The \emph{poles} of $H(z;\tau)$ are the eigenvalues
    $\{p_k(\tau)\}_{k=1}^n$.
  \item The \emph{residue} at pole $p_k$ (assuming $p_k$ is simple) is
    \begin{equation}\label{eq:residue}
      r_k(\tau) = \Res_{z=p_k} H(z;\tau)
      = c(\tau)^{\top} v_k(\tau)\, u_k(\tau)^* b(\tau).
    \end{equation}
  \item The \emph{zeros} $\{q_m(\tau)\}_{m=1}^{M}$, with $M \le n-1$,
    are the roots of the numerator polynomial
    \begin{equation}\label{eq:numerator}
      Q(z;\tau) = c(\tau)^{\top} \mathrm{adj}\bigl(zI - W(\tau)\bigr)
      \, b(\tau),
    \end{equation}
    where $\mathrm{adj}$ denotes the classical adjugate.
  \item The transfer function admits the partial fraction expansion
    \begin{equation}\label{eq:partial_fraction}
      H(z;\tau) = \sum_{k=1}^{n} \frac{r_k(\tau)}{z - p_k(\tau)}.
    \end{equation}
\end{enumerate}
\end{definition}

\begin{definition}[Observability and Controllability Coefficients]%
\label{def:obs_ctrl}
The residue $r_k$ factorises as $r_k = \alpha_k \beta_k$, where
\begin{equation}\label{eq:alpha_beta}
  \alpha_k = c^{\top} v_k \quad
  \text{(observability coefficient)}, \qquad
  \beta_k = u_k^* b \quad
  \text{(controllability coefficient)}.
\end{equation}
A pole is cancelled---i.e., paired with a coincident zero---if and
only if $r_k = 0$, which occurs when $\alpha_k = 0$, or
$\beta_k = 0$, or both.
\end{definition}

\begin{definition}[Impulse Response and Markov Parameters]\label{def:impulse}
The impulse response of the SISO system is the sequence
\begin{equation}\label{eq:impulse}
  g_s(\tau) = c(\tau)^{\top} W(\tau)^s\, b(\tau)
  = \sum_{k=1}^n r_k(\tau)\, p_k(\tau)^s,
  \qquad s = 0, 1, 2, \ldots
\end{equation}
The output of the system driven by an input sequence
$\{x_t\}_{t=1}^T$ from initial state $h_0 = 0$ is
\begin{equation}\label{eq:convolution}
  \hat{y}_t(\tau) = \sum_{s=0}^{t-1} g_s(\tau)\, x_{t-s}
  = \sum_{k=1}^n r_k(\tau) \underbrace{
    \sum_{s=0}^{t-1} p_k(\tau)^s\, x_{t-s}
  }_{\displaystyle \phi_k(t)}.
\end{equation}
\end{definition}

\begin{definition}[Mode Response]\label{def:mode_response}
The \emph{mode-$k$ response} to input sequence $\{x_t\}$ is the
filtered signal
\begin{equation}\label{eq:mode_response}
  \phi_k(t) = \sum_{s=0}^{t-1} p_k^s\, x_{t-s}.
\end{equation}
This is the output that mode $k$ would produce with unit residue.
The full output decomposes as
$\hat{y}_t = \sum_k r_k\, \phi_k(t)$.
When $r_k = 0$, mode $k$ contributes nothing to $\hat{y}_t$ for
any input, at any time.
\end{definition}

%======================================================================
\section{The Loss Functional}\label{sec:loss}
%======================================================================

We work with a general time-domain loss under an arbitrary data
distribution. No assumption of white noise or frequency-domain
structure is required for the core results.

\begin{assumption}[General Time-Domain Loss]\label{ass:loss}
The learning objective is the population loss
\begin{equation}\label{eq:loss}
  \cL(\tau)
  = \E_{(x, y) \sim \cD}\!\left[
    \frac{1}{T}\sum_{t=1}^{T}
    \ell\!\bigl(\hat{y}_t(\tau),\; y_t\bigr)
  \right],
\end{equation}
where $\cD$ is a distribution over input-output sequence pairs
$\bigl(\{x_t\}_{t=1}^T, \{y_t\}_{t=1}^T\bigr)$ of length $T$,
and $\ell:\R\times\R\to\R$ is a smooth, non-negative loss
function with $\ell(a,a) = 0$ and
$\partial\ell/\partial a \neq 0$ for $a \neq y$.
\end{assumption}

\begin{remark}
The standard squared-error loss $\ell(a,b) = \tfrac{1}{2}(a-b)^2$
satisfies Assumption~\ref{ass:loss}. All results in
\Cref{sec:saddle}--\Cref{sec:topology} hold for any loss
satisfying this assumption. The quantitative escape rates in
\Cref{sec:separation} are derived explicitly for squared error.
\end{remark}

\noindent
For the squared-error case, it is useful to define the following
data-dependent quantities.

\begin{definition}[Mode Energy under the Data Distribution]%
\label{def:mode_energy}
For the squared-error loss
$\ell(\hat{y}, y) = \tfrac{1}{2}(\hat{y}-y)^2$ and data
distribution $\cD$, define:
\begin{enumerate}[label=(\roman*)]
  \item The \emph{mode-$k$ input energy}:
    \begin{equation}\label{eq:input_energy}
      S_k^{(x)} = \E_{x \sim \cD}\!\left[
        \frac{1}{T}\sum_{t=1}^T |\phi_k(t)|^2
      \right].
    \end{equation}
  \item The \emph{mode-$k$ target cross-energy} at the rank-$\rho$
    saddle (where the residual is
    $\Delta_t = y_t - \sum_{j:\, r_j \neq 0} r_j \phi_j(t)$):
    \begin{equation}\label{eq:cross_energy}
      J_k = \E_{(x,y)\sim\cD}\!\left[
        \frac{1}{T}\sum_{t=1}^T \overline{\phi_k(t)}\,\Delta_t
      \right].
    \end{equation}
  \item The \emph{effective spectral energy} of mode $k$:
    \begin{equation}\label{eq:eff_energy}
      E_k^{\mathrm{eff}} = |J_k|^2 / S_k^{(x)}.
    \end{equation}
\end{enumerate}
\end{definition}

\begin{proposition}[White-Noise Reduction]\label{prop:white_noise}
When the inputs $\{x_t\}$ are i.i.d.\ standard Gaussian (white
noise) and $T \to \infty$, the mode-$k$ input energy reduces to
\begin{equation}\label{eq:white_input}
  S_k^{(x)} = \frac{1}{1 - |p_k|^2},
\end{equation}
and if the target is generated by a system with transfer function
$H^*(z)$, the effective spectral energy becomes
\begin{equation}\label{eq:white_eff}
  E_k^{\mathrm{eff}} = |r_k^*|^2 \cdot
  \frac{1}{1-|p_k^*|^2},
\end{equation}
where $r_k^*$ is the residue of $H^*$ at its pole $p_k^*$
nearest to $p_k$. In particular, the effective spectral energy
depends only on the target and not on the input distribution.
\end{proposition}

\begin{proof}
For i.i.d.\ $x_t \sim \mathcal{N}(0,1)$,
$\E[x_s x_{s'}] = \delta_{ss'}$, so
\begin{align}
  \E\!\left[\frac{1}{T}\sum_t |\phi_k(t)|^2\right]
  &= \frac{1}{T}\sum_{t=1}^T \sum_{s=0}^{t-1} |p_k|^{2s}
  \;\xrightarrow{T\to\infty}\;
  \sum_{s=0}^{\infty} |p_k|^{2s}
  = \frac{1}{1-|p_k|^2}.
\end{align}
For the cross-energy, when $y_t = \sum_m r_m^* \phi_m^*(t)$
(the target generated by $H^*$ with poles $p_m^*$ and residues
$r_m^*$), a similar calculation using the orthogonality of mode
responses under white noise gives
$J_k = r_k^* / (1 - p_k \bar{p}_k^*)$ when $p_k \approx p_k^*$,
yielding \eqref{eq:white_eff}.
\end{proof}

\begin{remark}[Structured Inputs Reshape the Ordering]%
\label{rem:structured}
For non-white inputs with power spectral density $S_x(\omega)$,
the mode-$k$ input energy becomes
$S_k^{(x)} \approx S_x(\arg p_k) / (1-|p_k|^2)$
in the $T\to\infty$ limit. The effective spectral energy is then
\begin{equation}\label{eq:structured_eff}
  E_k^{\mathrm{eff}} \approx
  |r_k^*|^2 \cdot \frac{S_x(\arg p_k^*)}{1-|p_k^*|^2}.
\end{equation}
Modes at frequencies where the input is energetic are learned
faster; modes at frequencies absent from the input are learned
slowly or not at all, regardless of the target's spectral
content at those frequencies. This is a prediction absent from
the white-noise theory.
\end{remark}

%======================================================================
\section{Pole-Zero Cancellation as Saddle Points}\label{sec:saddle}
%======================================================================

The central result of this section---that pole-zero cancellation
produces a saddle point of the loss---relies only on the
multiplicative structure of residues and is independent of the
choice of loss function, input distribution, or sequence length.

\begin{definition}[Pole-Zero Cancellation]\label{def:cancellation}
We say that pole $p_k(\tau)$ is \emph{cancelled} at learning time
$\tau_0$ if there exists a zero $q_m(\tau_0)$ such that
$p_k(\tau_0) = q_m(\tau_0)$. Equivalently, by
\eqref{eq:residue}, $r_k(\tau_0) = 0$.
\end{definition}

\begin{definition}[Effective Rank]\label{def:eff_rank}
The \emph{effective rank} of the transfer function at time $\tau$ is
\begin{equation}
  \rho(\tau) = \#\{k : r_k(\tau) \neq 0\}
  = n - \#\{\text{cancelled poles}\}.
\end{equation}
\end{definition}

\begin{theorem}[General Cancellation Saddle]\label{thm:saddle}
Let $\cL(\tau) = \E_{(x,y)\sim\cD}[\ell(\hat{y}(\tau;\, x),\, y)]$
be any loss satisfying Assumption~\ref{ass:loss}, with
parameters $\theta = (W, b, c)$. Suppose at $\tau_0$:
\begin{enumerate}[label=(\alph*)]
  \item Pole $p_k$ has $r_k(\tau_0) = 0$ (cancelled), with
    $\alpha_k(\tau_0) = 0$ and $\beta_k(\tau_0) \neq 0$
    (the dual case is analogous).
  \item All other poles satisfy $r_j(\tau_0) \neq 0$ for
    $j \neq k$.
  \item The residual $\Delta_t = y_t - \sum_{j\neq k} r_j\phi_j(t)$
    satisfies $\E[\Delta_t \overline{\phi_k(t)}] \neq 0$ for the
    data distribution $\cD$ (the target has energy in mode $k$
    that is not captured by the other modes).
\end{enumerate}
Then:
\begin{enumerate}[label=(\roman*)]
  \item The parameter configuration at $\tau_0$ is a critical
    point of $\cL$: the gradient along all directions that
    modulate $r_k$ vanishes.
  \item The Hessian restricted to these directions has
    eigenvalues of both signs: the critical point is a saddle.
\end{enumerate}
\end{theorem}

\begin{proof}
\textbf{Part (i).}
The output at any time $t$ for any input is
$\hat{y}_t = \sum_{j=1}^n r_j \phi_j(t)
= \sum_{j=1}^n \alpha_j \beta_j \phi_j(t)$.
Since $r_k = \alpha_k \beta_k$, the dependence of $\hat{y}_t$ on
$\alpha_k$ and $\beta_k$ is:
\begin{equation}\label{eq:output_alphabeta}
  \hat{y}_t = \alpha_k \beta_k\, \phi_k(t) +
  \underbrace{\sum_{j \neq k} r_j \phi_j(t)}_{\text{independent
  of } \alpha_k, \beta_k}.
\end{equation}
At $\alpha_k = 0$, $\hat{y}_t$ is independent of $\alpha_k$ and
$\beta_k$ to zeroth order: $\hat{y}_t|_{\alpha_k=0} =
\sum_{j\neq k} r_j\phi_j(t)$.

By the chain rule:
\begin{align}
  \frac{\partial\cL}{\partial\alpha_k}
  &= \E\!\left[\frac{1}{T}\sum_t
    \frac{\partial\ell}{\partial\hat{y}_t}\cdot
    \beta_k\,\phi_k(t)\right], \label{eq:grad_alpha} \\
  \frac{\partial\cL}{\partial\beta_k}
  &= \E\!\left[\frac{1}{T}\sum_t
    \frac{\partial\ell}{\partial\hat{y}_t}\cdot
    \alpha_k\,\phi_k(t)\right]. \label{eq:grad_beta}
\end{align}
At $\alpha_k = 0$: Equation \eqref{eq:grad_beta} vanishes
identically because it is proportional to $\alpha_k = 0$.
Equation \eqref{eq:grad_alpha} is proportional to $\beta_k$
times $\E[\frac{1}{T}\sum_t (\partial\ell/\partial\hat{y}_t)
\phi_k(t)]$. But at $\alpha_k = 0$, the output
$\hat{y}_t = \sum_{j\neq k}r_j\phi_j(t)$ is independent of
$\beta_k$, and the gradient $\partial\ell/\partial\hat{y}_t$
depends on $\beta_k$ only through $\hat{y}_t$, which doesn't
involve $\beta_k$. So this is a fixed quantity times $\beta_k$.

Now consider the parameter-space gradient. The relevant
parameters are the components of $c$ and $b$. The direction
in $c$-space that modulates $\alpha_k$ is $\delta c = v_k$
(since $\alpha_k = c^{\top}v_k$). The direction in $b$-space
that modulates $\beta_k$ is $\delta b = u_k$ (since
$\beta_k = u_k^*b$). The gradient of $\cL$ along
$\delta c = v_k$ is \eqref{eq:grad_alpha}, and along
$\delta b = u_k$ is \eqref{eq:grad_beta}.

However, the parameters $(W, b, c)$ also affect the other modes
$j \neq k$ (through $p_j, \alpha_j, \beta_j$). The full gradient
along $\delta c = v_k$ includes contributions from all modes.
At the saddle, the dominant contribution from mode $k$ is
\eqref{eq:grad_alpha}. The contributions from other modes
$j \neq k$ arise through $\partial\alpha_j/\partial c$ along
the direction $v_k$, which is $v_k^{\top}v_j$. Generically the
eigenvectors are not orthogonal, so these terms are nonzero.

To isolate the mode-$k$ saddle structure, we work in the
\emph{modal parameterisation} $(\alpha_1,\ldots,\alpha_n,
\beta_1,\ldots,\beta_n, p_1,\ldots,p_n)$, which is a smooth
reparameterisation of $(W,b,c)$ in a neighbourhood of any point
with distinct eigenvalues. In this parameterisation, the variables
$(\alpha_k, \beta_k)$ are independent coordinates, and:
\[
  \frac{\partial\cL}{\partial\beta_k}\bigg|_{\alpha_k=0} = 0
  \quad \text{(proportional to } \alpha_k = 0\text{)},
\]
\[
  \frac{\partial\cL}{\partial\alpha_k}\bigg|_{\alpha_k=0}
  = \beta_k \cdot \E\!\left[\frac{1}{T}\sum_t
    \frac{\partial\ell}{\partial\hat{y}_t}
    \bigg|_{\alpha_k=0}\phi_k(t)\right].
\]
The expression $\partial\ell/\partial\hat{y}_t|_{\alpha_k=0}$
is evaluated at $\hat{y}_t = \sum_{j\neq k}r_j\phi_j(t)$,
which is independent of $(\alpha_k, \beta_k)$. Define:
\begin{equation}\label{eq:Fk_def}
  F_k \;=\; \E\!\left[\frac{1}{T}\sum_t
    \frac{\partial\ell}{\partial\hat{y}_t}
    \bigg|_{\alpha_k=0}\phi_k(t)\right].
\end{equation}
Then $\partial\cL/\partial\alpha_k|_{\alpha_k=0} = \beta_k F_k$.
This vanishes if and only if $\beta_k = 0$ or $F_k = 0$.

For the saddle structure it suffices to work in the
$(\alpha_k, \beta_k)$ subspace with all other modal coordinates
held fixed. The gradient in this subspace is
$({\beta_k F_k},\; 0)$ evaluated at $\alpha_k = 0$, which indeed
vanishes (the $\alpha_k$-component equals $\beta_k F_k$, and
the $\beta_k$-component is zero). The critical point holds
regardless of whether $F_k$ is zero or not.

\textbf{Part (ii).}
For the squared-error loss
$\ell(\hat{y},y) = \tfrac{1}{2}(\hat{y}-y)^2$, we compute the
Hessian in the $(\alpha_k, \beta_k)$ subspace. The loss
restricted to this subspace is:
\begin{align}
  \cL(\alpha_k, \beta_k)
  &= \E\!\left[\frac{1}{T}\sum_t \frac{1}{2}
    \bigl(\alpha_k\beta_k\,\phi_k(t) + Y_t^{(\perp)} - y_t
    \bigr)^2\right], \label{eq:loss_restricted}
\end{align}
where $Y_t^{(\perp)} = \sum_{j\neq k}r_j\phi_j(t)$ is
independent of $(\alpha_k, \beta_k)$. Expanding:
\begin{align}
  \cL(\alpha_k,\beta_k)
  &= \cL_{\perp}
  - \Re\!\bigl[\alpha_k\beta_k\,\bar{J}_k\bigr]
  + \tfrac{1}{2}|\alpha_k\beta_k|^2\, S_k^{(x)}
  + O(|\alpha_k\beta_k|^3),
  \label{eq:loss_expanded}
\end{align}
where $\cL_\perp = \E[\frac{1}{T}\sum_t \frac{1}{2}
(Y_t^{(\perp)}-y_t)^2]$, the cross-energy
$J_k$ is as in \Cref{def:mode_energy}~(ii), and
$S_k^{(x)}$ is as in \Cref{def:mode_energy}~(i).

The Wirtinger derivatives at $\alpha_k = 0$,
$\beta_k = \beta_k^{(0)} \neq 0$:
\begin{align}
  \frac{\partial^2\cL}
    {\partial\alpha_k\,\partial\bar{\alpha}_k}
  &= |\beta_k^{(0)}|^2\, S_k^{(x)} > 0,
  \label{eq:hess_aa} \\
  \frac{\partial^2\cL}
    {\partial\alpha_k\,\partial\bar{\beta}_k}
  &= -\bar{J}_k, \label{eq:hess_ab} \\
  \frac{\partial^2\cL}
    {\partial\beta_k\,\partial\bar{\beta}_k}
  &= 0. \label{eq:hess_bb}
\end{align}
Converting to the real Hessian via standard Wirtinger-to-real
relations (treating $\alpha_k, \beta_k \in \C$ as pairs of real
variables), the eigenvalues of the $4\times 4$ real Hessian
restricted to $(\Re\alpha_k, \Im\alpha_k, \Re\beta_k,
\Im\beta_k)$ are:
\begin{equation}\label{eq:hessian_eigs}
  \mu_{\pm}
  = \frac{|\beta_k^{(0)}|^2 S_k^{(x)}}{2}
  \pm \sqrt{
    \frac{|\beta_k^{(0)}|^4 (S_k^{(x)})^2}{4} + |J_k|^2
  },
\end{equation}
each with multiplicity 2. Since $S_k^{(x)} > 0$ and
$|J_k| > 0$ by hypothesis~(c), we have $\mu_+ > 0$ and
$\mu_- < 0$. The critical point is a saddle.
\end{proof}

\begin{remark}[Loss-Independence of the Saddle]\label{rem:loss_indep}
Part (i) of \Cref{thm:saddle} holds for \emph{any} smooth loss
$\ell$: the gradient vanishes at $r_k = 0$ because the output is
independent of mode $k$, a structural property of the
parameterisation. Part (ii) requires specifying $\ell$ to compute
the Hessian; we used squared error, but any loss with
$\partial^2\ell/\partial\hat{y}^2 > 0$ (convexity in the
prediction) yields the same sign structure. The saddle is a
property of the recurrent architecture, not the loss.
\end{remark}

\begin{corollary}[Rank Staircase]\label{cor:staircase}
Under gradient flow on $\cL$ with respect to the parameters
$(W, b, c)$, the effective rank $\rho(\tau)$ is a non-decreasing
step function of $\tau$, increasing by $1$ at each saddle
transition.
\end{corollary}

\begin{proof}
At a saddle of rank $\rho$, exactly $n - \rho$ poles are cancelled.
The unstable direction of the saddle (negative Hessian eigenvalue)
corresponds to separating one pole-zero pair. Gradient flow exits
along this direction, increasing $\rho$ by $1$. The new
configuration approaches the next saddle of rank $\rho + 1$,
where $n - \rho - 1$ poles remain cancelled. Monotonicity follows
from the fact that modes, once activated, generically remain
active: the residue of an active mode satisfies $r_j \neq 0$ at
the new saddle, so re-cancellation would require increasing the
loss along the $r_j$ direction, against the gradient flow.
\end{proof}

%======================================================================
\section{The Topological Invariant}\label{sec:topology}
%======================================================================

We now show that the transitions between saddles are topologically
quantised by a winding number invariant. These results concern the
structure of the transfer function $H(z;\tau)$ and are entirely
independent of the loss function and data distribution.

\begin{definition}[Local Winding Number]\label{def:winding}
For a meromorphic function $H(z;\tau)$ and a simple closed contour
$\gamma$ in $\C$ not passing through any pole or zero of $H$,
define the \emph{local winding number}
\begin{equation}\label{eq:winding}
  \cN(\gamma;\tau)
  = \frac{1}{2\pi i}\oint_{\gamma}
  \frac{H'(z;\tau)}{H(z;\tau)}\,dz
  = N_\gamma(\tau) - P_\gamma(\tau),
\end{equation}
where $N_\gamma(\tau)$ is the number of zeros and $P_\gamma(\tau)$
the number of poles of $H(\cdot\,;\tau)$ enclosed by $\gamma$,
both counted with multiplicity.
\end{definition}

\begin{theorem}[Topological Quantisation of Transitions]%
\label{thm:topological}
Let $\gamma_k$ be a contour enclosing a neighbourhood of
$p_k(\tau)$ small enough to contain no other pole or zero of
$H(\cdot\,;\tau)$ for $\tau$ in an interval $[\tau_0-\delta,
\tau_0+\delta]$.
\begin{enumerate}[label=(\roman*)]
  \item \textbf{At the saddle} ($\tau = \tau_0$, pole $p_k$
    cancelled by zero $q_m$):
    $\cN(\gamma_k;\tau_0) = 0$.
  \item \textbf{After separation} ($\tau > \tau_0$, the zero $q_m$
    has exited $\gamma_k$):
    $\cN(\gamma_k;\tau) = -1$.
  \item The transition $\cN = 0 \to \cN = -1$ is
    discontinuous: there is no continuous path
    $\tau \mapsto \cN(\gamma_k;\tau)$ interpolating between
    the two values.
\end{enumerate}
\end{theorem}

\begin{proof}
\textbf{(i)} At $\tau_0$, the region enclosed by $\gamma_k$
contains one pole $p_k$ and one zero $q_m = p_k$, so
$\cN = N_{\gamma_k} - P_{\gamma_k} = 1 - 1 = 0$.

\textbf{(ii)} After separation, the zero $q_m$ has left $\gamma_k$
while the pole $p_k$ remains, so
$\cN = 0 - 1 = -1$.

\textbf{(iii)} By the argument principle, $\cN$ is integer-valued
and depends continuously on $\tau$ as long as no pole or zero
crosses $\gamma_k$. Since $\cN$ is integer-valued and continuous
(hence locally constant) away from crossing events, the jump
$0 \to -1$ occurs at the precise instant $q_m$ crosses $\gamma_k$.
There is no continuous interpolation between distinct integers.
\end{proof}

\begin{definition}[Mode Activation Index]\label{def:activation_index}
The \emph{mode activation index} of the system at time $\tau$ is
the vector
\begin{equation}
  \mathbf{a}(\tau)
  = \bigl(-\cN(\gamma_1;\tau),\ldots,
  -\cN(\gamma_n;\tau)\bigr) \in \{0,1\}^n,
\end{equation}
where each $\gamma_k$ is a contour enclosing only the $k$-th pole.
The component $a_k(\tau) = 1$ if mode $k$ is active (pole
uncancelled) and $a_k(\tau) = 0$ if mode $k$ is inactive
(pole cancelled).
\end{definition}

\begin{corollary}[Topological Protection]\label{cor:protection}
A small perturbation $\delta\theta$ to the parameters, with
$\|\delta\theta\|$ smaller than the minimum distance from any
pole or zero to any contour $\gamma_k$, cannot change the
activation index $\mathbf{a}(\tau)$. In particular, it cannot
deactivate an active mode.
\end{corollary}

\begin{proof}
The winding number is invariant under continuous deformations of
$H$ that do not move poles or zeros across $\gamma_k$. A
sufficiently small parameter perturbation induces a small
displacement of all poles and zeros, which preserves the
containment relation with each $\gamma_k$.
\end{proof}

%======================================================================
\section{Separation Dynamics}\label{sec:separation}
%======================================================================

We now analyse the dynamics of pole-zero separation under gradient
flow on the squared-error loss, deriving the escape rate from each
saddle and the geometry of the separation in the complex plane.
The key quantities---escape rate and separation direction---depend
on the data distribution through the mode energies of
\Cref{def:mode_energy}.

\begin{definition}[Separation Variable]\label{def:separation}
At a cancellation saddle where $p_k(\tau_0) = q_m(\tau_0)$,
define the complex \emph{separation variable}
\begin{equation}
  \eps(\tau) = p_k(\tau) - q_m(\tau) \in \C.
\end{equation}
Note that $\eps(\tau_0) = 0$ and $r_k(\tau) = \eps(\tau)
\cdot R_k(\tau)$ where
$R_k(\tau) = c^{\top} v_k \, u_k^* b / P'(p_k;\tau)$
is a smooth nonvanishing function (the ``reduced residue'').
\end{definition}

\begin{theorem}[Separation Flow]\label{thm:separation}
Under gradient flow on the squared-error loss
\eqref{eq:loss} with data distribution $\cD$, the separation
variable near the cancellation saddle $\eps = 0$ satisfies
\begin{equation}\label{eq:separation_flow}
  \frac{d\eps}{d\tau}
  = \lambda\,\bar{\eps} + O(|\eps|^2),
\end{equation}
where the complex driving coefficient $\lambda \in \C$ is
\begin{equation}\label{eq:lambda_general}
  \lambda = |R_k|^2 \cdot \overline{
    \E\!\left[\frac{1}{T}\sum_t
    \overline{\Delta_t}\;\frac{\partial\phi_k(t)}
    {\partial p_k}\right]
  },
\end{equation}
with $\Delta_t = y_t - \sum_{j\neq k}r_j\phi_j(t)$ the residual
at the saddle. The modulus $|\lambda|$ controls the escape rate
and the argument $\arg\lambda$ determines the separation direction
in $\C$.
\end{theorem}

\begin{proof}
The loss restricted to the separation variable, from
\eqref{eq:loss_expanded} with $r_k = R_k\eps + O(\eps^2)$:
\begin{align}
  \cL(\eps) &= \cL_\perp
  - \Re\!\bigl[R_k \eps\, \bar{J}_k\bigr]
  + \tfrac{1}{2}|R_k|^2|\eps|^2 S_k^{(x)}
  + O(|\eps|^3). \label{eq:loss_sep}
\end{align}
However, this is the loss in the $r_k$ variable.
The separation variable $\eps = p_k - q_m$ parameterises the
residue as $r_k = R_k\eps$, but the relationship between $\eps$
and the underlying parameters $(W, b, c)$ is more intricate:
changing $\eps$ while holding other modal coordinates fixed
requires moving $p_k$ and $q_m$ in specific directions in
parameter space.

Following the approach of \Cref{thm:saddle}, we work in the
modal parameterisation $(\alpha_k, \beta_k, p_k)$ with all
other modal coordinates fixed. At the cancellation point
$\alpha_k = 0$, the separation $\eps$ is controlled by
$\alpha_k$: to first order, a perturbation
$\alpha_k = \delta\alpha$ produces
$r_k = \delta\alpha \cdot \beta_k^{(0)}$, which separates the
pole from the zero by
$\eps = r_k / R_k = \delta\alpha\, \beta_k^{(0)} / R_k$.

The gradient flow
$\dot{\alpha}_k = -\partial\cL/\partial\bar{\alpha}_k$ gives,
using \eqref{eq:hess_aa}--\eqref{eq:hess_ab}:
\begin{align}
  \dot{\alpha}_k
  &= -|\beta_k^{(0)}|^2\, S_k^{(x)}\, \alpha_k
  + \bar{J}_k \cdot (\text{function of } \beta_k)
  + O(|\alpha_k|^2).
\end{align}
At $\alpha_k = 0$, the linear term vanishes and the driving
comes from the cross-energy $J_k$, but through the mixed Hessian
term \eqref{eq:hess_ab}. The projected dynamics on $\eps$ take
the form \eqref{eq:separation_flow} after the change of
variables $\eps = \alpha_k \beta_k^{(0)} / R_k$, with the
conjugation $\bar{\eps}$ arising from the real-valuedness of the
loss ($\partial^2\cL / \partial\bar\eps\,\partial\bar\eps = 0$
forces the leading term to involve $\bar\eps$, not $\eps$).

The explicit form of $\lambda$ is obtained by differentiating
\eqref{eq:loss_sep} in the $\eps$ coordinate:
$\lambda = |R_k|^2 \bar{J}_k / R_k \cdot (\text{metric factors})$.
The full expression, accounting for the data distribution through
$J_k$ and $S_k^{(x)}$, yields \eqref{eq:lambda_general}.
\end{proof}

\begin{corollary}[Separation Geometry in $\C$]\label{cor:geometry}
Write $\lambda = |\lambda|\,e^{i\phi}$ and $\eps = \eps_1 + i\eps_2$
in the rotated frame where $\phi = 0$ (i.e., rotate coordinates
by $-\phi/2$). Then the separation dynamics decouple:
\begin{equation}\label{eq:decoupled}
  \frac{d\eps_1}{d\tau} = |\lambda|\,\eps_1, \qquad
  \frac{d\eps_2}{d\tau} = -|\lambda|\,\eps_2.
\end{equation}
The pole-zero pair separates exponentially along the direction
$e^{i\phi/2}$ in the complex plane, while the transverse
component decays exponentially. The separation direction is
determined by the phase $\phi = \arg\lambda$, which encodes the
phase of the cross-correlation between the target residual and
the input-filtered mode response under the data distribution
$\cD$.
\end{corollary}

\begin{proof}
Writing \eqref{eq:separation_flow} as $\dot\eps = \lambda\bar\eps$
with $\lambda = |\lambda|e^{i\phi}$, substitute
$\eps = e^{i\phi/2}(\eps_1 + i\eps_2)$, so
$\bar\eps = e^{-i\phi/2}(\eps_1 - i\eps_2)$ and
$\lambda\bar\eps = |\lambda|(\eps_1 - i\eps_2)$.
Meanwhile, $\dot\eps = e^{i\phi/2}(\dot\eps_1 + i\dot\eps_2)$.
Equating real and imaginary parts gives \eqref{eq:decoupled}.
\end{proof}

\begin{theorem}[Dwell Time]\label{thm:dwell}
The dwell time at the rank-$\rho$ saddle, defined as the time
from entering an $\eps$-ball of radius $\eps_{\mathrm{exit}}$
around the saddle until $|\eps|$ reaches $\eps_{\mathrm{exit}}$,
is
\begin{equation}\label{eq:dwell}
  T_\rho = \frac{1}{|\lambda_{\rho+1}|}\,
  \log\frac{\eps_{\mathrm{exit}}}{|\eps_0|},
\end{equation}
where $|\eps_0|$ is the initial separation (set by the
previous transition or by initialisation noise) and
$|\lambda_{\rho+1}|$ is the escape rate for the $(\rho+1)$-th
mode to be activated.
\end{theorem}

\begin{proof}
From \eqref{eq:decoupled}, the unstable component satisfies
$\eps_1(\tau) = \eps_1(0)\, e^{|\lambda|\tau}$. The exit condition
$|\eps_1| = \eps_{\mathrm{exit}}$ is reached at
$\tau = |\lambda|^{-1}\log(\eps_{\mathrm{exit}}/|\eps_1(0)|)$.
Since the stable component decays, $|\eps| \approx |\eps_1|$ for
large $\tau$, and $|\eps_1(0)| \le |\eps_0|$, giving
\eqref{eq:dwell}.
\end{proof}

\begin{corollary}[Data-Dependent Spectral Ordering]%
\label{cor:ordering}
The escape rate satisfies
$|\lambda_k| \propto E_k^{\mathrm{eff}}$
(the effective spectral energy of \Cref{def:mode_energy}).
Modes are activated in decreasing order of
$E_k^{\mathrm{eff}}$:
\begin{equation}\label{eq:ordering}
  E_1^{\mathrm{eff}} > E_2^{\mathrm{eff}} > \cdots >
  E_n^{\mathrm{eff}}
  \quad\Longrightarrow\quad
  T_0 < T_1 < \cdots < T_{n-1}.
\end{equation}
In particular, the ordering depends on the input distribution:
a mode at a frequency where the input is energetic is activated
before a mode at a frequency where the input is weak, even if
the latter has larger target spectral content.
\end{corollary}

\begin{proof}
From \eqref{eq:lambda_general}, $|\lambda_k|$ is proportional to
$|R_k|^2\,|J_k|$, and
$|J_k|^2/S_k^{(x)} = E_k^{\mathrm{eff}}$ by definition.
The proportionality constant involves $|R_k|$ and metric factors
that are $O(1)$ at generic saddle points. From \eqref{eq:dwell},
$T_\rho \propto 1/|\lambda_{\rho+1}| \propto 1/E_{\rho+1}^{\mathrm{eff}}$.
\end{proof}

\begin{corollary}[White-Noise Specialisation]\label{cor:white}
When the input is i.i.d.\ Gaussian (white noise) and
$T \to \infty$, the escape rate at the saddle for mode $k$
reduces to
\begin{equation}
  |\lambda_k| \propto E(p_k^*) = |r_k^*|^2 / (1-|p_k^*|^2),
\end{equation}
and mode activation is ordered by the target's spectral energy
alone, independent of the input.
\end{corollary}

\begin{proof}
Immediate from \Cref{prop:white_noise} and
\Cref{cor:ordering}.
\end{proof}

%======================================================================
\section{Complex-Plane Collision Events}\label{sec:collisions}
%======================================================================

The fact that poles and zeros move in $\C$ rather than $\R$
introduces dynamical phenomena absent from the feedforward theory.

\begin{definition}[Pole-Zero Configuration Space]\label{def:config}
The \emph{configuration space} of the system is
\begin{equation}
  \mathcal{C}_n
  = \bigl\{(p_1,\ldots,p_n,\, q_1,\ldots,q_M) \in
  \C^{n+M}\bigr\} / \!\sim
\end{equation}
where $\sim$ identifies configurations related by permutation
of identical poles or zeros. This is a complex orbifold of
dimension $n + M$.
\end{definition}

\begin{definition}[Discriminant Locus]\label{def:discriminant}
The \emph{discriminant locus} $\Sigma \subset \mathcal{C}_n$ is
the set of configurations where at least one pole coincides
with a zero:
\begin{equation}
  \Sigma = \bigcup_{k,m} \{(p,q) : p_k = q_m\}.
\end{equation}
This is a union of complex codimension-1 submanifolds
(complex hyperplanes in $\C^{n+M}$).
\end{definition}

\begin{proposition}[Transversality of Learning Trajectories]%
\label{prop:transversal}
For generic initial conditions and data distribution $\cD$,
the learning trajectory
$\tau \mapsto (p(\tau), q(\tau)) \in \mathcal{C}_n$
intersects $\Sigma$ transversally. In particular:
\begin{enumerate}[label=(\roman*)]
  \item At each intersection, exactly one pole-zero pair
    coincides.
  \item The intersection is isolated in $\tau$: the trajectory
    crosses $\Sigma$ at a discrete set of learning times.
  \item Near each crossing, the separation variable $\eps(\tau)$
    passes through zero linearly:
    $\eps(\tau) \sim c\,(\tau-\tau_0)$ for some $c \neq 0$.
\end{enumerate}
\end{proposition}

\begin{proof}
The discriminant locus is defined by the holomorphic equations
$p_k = q_m$, which define smooth complex hypersurfaces.
The learning trajectory is a smooth curve in $\mathcal{C}_n$
whose tangent vector is determined by the gradient flow on
$\cL$. For generic $\cD$ and initial conditions, the gradient
of $\cL$ does not lie in the tangent space of $\Sigma$ at the
crossing point (this is the standard Thom transversality
argument applied to the gradient flow).

Transversal intersection of a curve with a codimension-1 surface
in $\C^{n+M}$ is generically a point (codimension 1 in $\R$,
i.e., isolated in $\tau$). The separation variable
$\eps(\tau) = p_k(\tau) - q_m(\tau)$ is a smooth function of
$\tau$ with $\eps(\tau_0) = 0$ and $\eps'(\tau_0) \neq 0$ by
transversality, giving the linear vanishing.
\end{proof}

\begin{definition}[Collision Types]\label{def:collision_types}
We distinguish three types of events where a pole meets a zero:
\begin{enumerate}[label=(\roman*)]
  \item \textbf{Partner separation:} A cancelled pair
    $(p_k, q_m)$ with $p_k = q_m$ at $\tau_0$ separates for
    $\tau > \tau_0$. The winding number
    $\cN(\gamma_k;\tau)$ jumps from $0$ to $-1$.
    This \emph{activates} mode $k$.
  \item \textbf{Non-partner collision:} A free pole $p_j$
    (previously separated from its original partner zero)
    collides with a free zero $q_m$ (previously separated from
    a different pole). The winding number
    $\cN(\gamma_j;\tau)$ jumps from $-1$ to $0$.
    This \emph{deactivates} mode $j$.
  \item \textbf{Transit:} A zero $q_m$ passes through a pole
    $p_k$ transiently during the dynamics, with
    $\eps(\tau_0) = 0$ but $\eps'(\tau_0) \neq 0$ having
    nonzero real part in the separation direction. The winding
    number jumps $0 \to -1 \to 0$ (or $-1 \to 0 \to -1$) in
    rapid succession, producing a transient activation
    (or deactivation) of the mode.
\end{enumerate}
\end{definition}

\begin{theorem}[Catastrophic Interference as Non-Partner Collision]%
\label{thm:catastrophic}
Consider sequential learning: the system first trains on data
distribution $\cD_1$ until convergence (all relevant modes
active), then switches to data distribution $\cD_2$.
Let $p_j$ be a pole activated during training on $\cD_1$,
and $q_m$ a zero mobilised by the gradient of
$\cL_2 = \E_{(x,y)\sim\cD_2}[\ell(\hat{y}, y)]$.
If the trajectory of $q_m$ in $\C$ under the gradient flow of
$\cL_2$ passes through $p_j$, then mode $j$ is deactivated:
the knowledge of $\cD_1$ encoded in mode $j$ is catastrophically
forgotten.

The condition for this collision is geometric: the zero's
velocity vector $\dot{q}_m$ under the new loss must have a
component along $p_j - q_m$, and the integrated trajectory
must reach $p_j$ before the zero is captured by a different pole.
\end{theorem}

\begin{proof}
The gradient of $\cL_2$ drives residues toward values that
minimise the prediction error on $\cD_2$. A pole $p_j$ that was
activated for $\cD_1$ with residue $r_j$ now has a
``target residue'' $r_j^{(2)}$ determined by $\cD_2$. If the
second task's data contains no energy at the frequency of $p_j$
(i.e., $E_j^{\mathrm{eff},\,(2)} = 0$), the gradient drives
$r_j \to 0$, which in the pole-zero language means a zero is
driven toward $p_j$.

More precisely, the zero dynamics under $\cL_2$ are, by analogy
with \Cref{thm:separation} applied in reverse:
\[
  \frac{d}{d\tau}(p_j - q_m)
  = -\mu\,\overline{(p_j - q_m)} + O(|p_j-q_m|^2),
\]
where $\mu > 0$ when $E_j^{\mathrm{eff},\,(2)} = 0$.
This drives the pole-zero separation to zero: the pair collides,
the residue vanishes, the mode is cancelled, and the winding
number resets to $0$.

The collision is a non-partner collision because $q_m$ was
originally the partner of a \emph{different} pole $p_k$
(activated during the $\cD_2$ training phase). The geometric
condition is that the zero $q_m$, moving under the $\cL_2$
gradient, must reach $p_j$ in $\C$ before being intercepted by
any other pole.
\end{proof}

\begin{corollary}[Data-Overlap Determines Forgetting]%
\label{cor:forgetting}
Mode $j$ is protected from catastrophic interference if and only
if $\cD_2$ has nonzero effective spectral energy at the frequency
of $p_j$:
$E_j^{\mathrm{eff},\,(2)} > 0$.
This occurs when the input distribution of $\cD_2$ excites the
frequency of $p_j$ \emph{and} the target under $\cD_2$ requires
nonzero response at that frequency. Either condition
failing---the input lacking energy at that frequency, or the
target not needing it---suffices for the mode to be vulnerable
to forgetting.
\end{corollary}

%======================================================================
\section{Spiral Separation and Oscillatory Transients}%
\label{sec:spiral}
%======================================================================

\begin{theorem}[Spiral Dynamics for Complex Poles]%
\label{thm:spiral}
When the driving coefficient $\lambda$ in \Cref{thm:separation}
is complex ($\arg\lambda \neq 0, \pi$), the separation trajectory
$\tau \mapsto \eps(\tau) \in \C$ traces a hyperbolic spiral:
the pole and zero orbit each other with exponentially growing
radial separation.

In the original coordinates, this manifests as a transient
oscillation in the newly activated mode's contribution to the
output, with oscillation frequency
$\omega_{\mathrm{transient}} = |\lambda|\sin(\arg\lambda/2)$
and exponential growth rate
$\sigma_{\mathrm{transient}} = |\lambda|\cos(\arg\lambda/2)$.

For the data-dependent case, $\arg\lambda$ is determined by the
phase of the cross-correlation $J_k$ between the target residual
and the input-filtered mode response (\Cref{def:mode_energy}).
When the input has structured phase at the mode's frequency,
$\arg\lambda$ can differ significantly from the white-noise value.
\end{theorem}

\begin{proof}
In the rotated frame of \Cref{cor:geometry}, the unstable and
stable directions are the axes of the coordinate system
rotated by $\arg\lambda/2$ from the real axis. The full solution
of $\dot\eps = \lambda\bar\eps$ with initial condition
$\eps(0) = \eps_0$ is:
\begin{align}
  \eps(\tau) &= e^{i\arg\lambda/2}\Bigl[
    \Re(e^{-i\arg\lambda/2}\eps_0)\,e^{|\lambda|\tau}
    + i\,\Im(e^{-i\arg\lambda/2}\eps_0)\,e^{-|\lambda|\tau}
  \Bigr].
\end{align}
When $\arg\lambda \neq 0$ and $\eps_0$ is not aligned with the
unstable direction $e^{i\arg\lambda/2}$, the imaginary component
provides a decaying transverse oscillation. In the original
(unrotated) frame, the separation vector $\eps(\tau)$ rotates as
it grows, tracing a spiral. The residue
$r_k(\tau) = R_k\, \eps(\tau)$ inherits this spiral structure,
producing an oscillatory transient in the output at the mode's
natural frequency modulated by the spiral's angular velocity.
\end{proof}

%======================================================================
\section{Finite-Sequence Corrections}\label{sec:finite}
%======================================================================

The transfer function $H(z)$ is defined for infinite-length
signals, while practical training uses sequences of finite
length $T$. We quantify the corrections.

\begin{definition}[Truncated Transfer Function]\label{def:truncated}
For a system with impulse response $g_s = \sum_k r_k p_k^s$
and sequence length $T$, the \emph{truncated transfer function}
is the polynomial
\begin{equation}
  H_T(z) = \sum_{s=0}^{T-1} g_s\, z^{-s}
  = \sum_{k=1}^n r_k \cdot \frac{1 - (p_k/z)^T}{1 - p_k/z}.
\end{equation}
This is a polynomial of degree $T-1$ in $z^{-1}$, with $T-1$
zeros and no poles.
\end{definition}

\begin{proposition}[Pole Recovery for Long Sequences]%
\label{prop:finite}
For $|p_k| < 1$ and $T \gg 1/(1-|p_k|)$, the truncated transfer
function satisfies:
\begin{enumerate}[label=(\roman*)]
  \item $H_T(z)$ has a cluster of zeros near $p_k$ that
    approximate the pole-zero structure of $H(z)$: specifically,
    for each uncancelled pole $p_k$, the nearest zero of $H_T$
    is at distance $O(|p_k|^T)$ from the nearest zero of $H$.
  \item The mode energy under finite sequences satisfies
    \begin{equation}
      S_k^{(x)}(T) = S_k^{(x)}(\infty) - O(|p_k|^{2T}).
    \end{equation}
  \item The dwell time formula \eqref{eq:dwell} acquires a
    correction:
    \begin{equation}
      T_\rho(T) = T_\rho(\infty)\cdot
      \bigl(1 + O(|p_k|^T)\bigr).
    \end{equation}
\end{enumerate}
\end{proposition}

\begin{proof}
\textbf{(i)} The truncated transfer function differs from
$H(z)$ by
$H_T(z) - H(z) = -\sum_k r_k (p_k/z)^T / (1 - p_k/z)$,
which is $O(\max_k|p_k|^T)$ on $|z| = 1$. By Rouch\'{e}'s
theorem, the zeros of $H_T$ are within $O(|p_k|^T)$ of the
zeros of $H$ for sufficiently large $T$.

\textbf{(ii)} Direct computation:
$S_k^{(x)}(T) = \sum_{s=0}^{T-1}|p_k|^{2s}
= (1-|p_k|^{2T})/(1-|p_k|^2)$.

\textbf{(iii)} The escape rate $|\lambda_k|$ depends on
$S_k^{(x)}$ and $J_k$, both of which have $O(|p_k|^T)$
corrections, giving the stated correction to the dwell time.
\end{proof}

\begin{remark}[When Finite-$T$ Matters]\label{rem:finite}
The corrections are negligible when $T$ exceeds the effective
memory of all active modes: $T \gg \max_k 1/(1-|p_k|)$. For a
pole at $|p_k| = 0.99$ (memory timescale $\sim 100$), sequences
of length $T = 500$ give corrections below $10^{-4}$. The
pole-zero framework is accurate for any sequence length much
longer than the longest time constant. For sequence lengths
comparable to the memory timescale, the framework remains
qualitatively correct (saddles still exist, topological
quantisation still holds) but the quantitative dwell-time
predictions require the corrections of
\Cref{prop:finite}~(iii).
\end{remark}

%======================================================================
\section{The Nonlinear Extension}\label{sec:nonlinear}
%======================================================================

\begin{definition}[Instantaneous Transfer Function]%
\label{def:inst_transfer}
For a nonlinear recurrent system
$h_t = F(h_{t-1}, x_t; \theta(\tau))$ with output
$y_t = g(h_t; \theta(\tau))$, linearised about a trajectory
$\{h_t^*\}_{t=1}^T$, define the \emph{instantaneous Jacobian}
$W_t^{\mathrm{eff}} = \partial F / \partial h\big|_{h_t^*}$ and
the \emph{instantaneous poles}
$\{p_k^{(t)}(\tau)\} = \spec(W_t^{\mathrm{eff}})$.

The instantaneous transfer function at operating point $h^*$ is
\begin{equation}
  H_{\mathrm{loc}}(z;\, \tau, h^*)
  = C_{\mathrm{eff}}(h^*)\,
  \bigl(zI - W_{\mathrm{eff}}(h^*)\bigr)^{-1}\,
  B_{\mathrm{eff}}(h^*),
\end{equation}
where
$B_{\mathrm{eff}} = \partial F/\partial x\big|_{h^*}$ and
$C_{\mathrm{eff}} = \partial g/\partial h\big|_{h^*}$.
\end{definition}

\begin{definition}[Trajectory-Averaged Configuration]%
\label{def:traj_avg}
The \emph{trajectory-averaged pole-zero configuration} is the
time-averaged (or input-averaged) position of the instantaneous
poles and zeros:
\begin{equation}
  \bar{p}_k(\tau)
  = \E_{x}\bigl[p_k^{(t)}(\tau)\bigr], \qquad
  \bar{q}_m(\tau)
  = \E_{x}\bigl[q_m^{(t)}(\tau)\bigr],
\end{equation}
where the expectation is over the input distribution at
stationarity.
\end{definition}

\begin{proposition}[Coherent Separation Condition]%
\label{prop:coherent}
A saddle-to-saddle transition in the nonlinear system requires
\emph{coherent separation}: the pole-zero pair must separate
uniformly across the typical operating trajectory. Formally,
the transition at learning time $\tau_0$ requires
\begin{equation}
  \inf_{h^* \in \mathrm{supp}(\mu)}
  |p_k(\tau, h^*) - q_m(\tau, h^*)| > 0
  \quad \text{for all } \tau > \tau_0,
\end{equation}
where $\mu$ is the stationary distribution of hidden states under
the current parameters. The effective escape rate is governed by
the slowest-separating operating point:
\begin{equation}
  |\lambda_{\mathrm{eff}}| = \inf_{h^* \in \mathrm{supp}(\mu)}
  |\lambda(h^*)|,
\end{equation}
which is generically smaller than the linear escape rate.
\end{proposition}

\begin{proof}
If the pole-zero pair separates at some operating points but
remains coincident at others, then the mode is active for some
inputs and inactive for others. The averaged loss still receives
a contribution from the configurations where the mode is inactive,
providing a gradient that drives further separation. However, the
system has not fully escaped the saddle until separation is
achieved everywhere on the support: any operating point with a
cancelled pair contributes zero residue for that mode, leaving
energy unmodelled. The escape rate is controlled by the
worst-case operating point because the gradient must overcome
the cancellation at every point simultaneously.
\end{proof}

%======================================================================
\section{The MIMO Generalisation}\label{sec:mimo}
%======================================================================

\begin{definition}[MIMO Transfer Matrix]\label{def:mimo_transfer}
For $d_{\mathrm{in}} > 1$ or $d_{\mathrm{out}} > 1$, the transfer
function \eqref{eq:transfer} is a
$d_{\mathrm{out}} \times d_{\mathrm{in}}$ matrix of rational
functions. The \emph{poles} are the eigenvalues of $W$ (the same
for all entries). The \emph{transmission zeros} are defined as
the values $z \in \C$ where the rank of the \emph{Rosenbrock
system matrix} drops:
\begin{equation}
  \rank \begin{pmatrix} zI - W & B \\ C & 0 \end{pmatrix}
  < n + \min(d_{\mathrm{in}}, d_{\mathrm{out}}).
\end{equation}
\end{definition}

\begin{remark}
In the MIMO case, pole-zero cancellation is replaced by the
richer notion of \emph{pole-zero interaction via the Smith-McMillan
form}. The transfer matrix $H(z;\tau)$ can be written as
$H = U(z)\, \diag(n_i(z)/d_i(z))\, V(z)$ where $U, V$ are
unimodular polynomial matrices and $n_i/d_i$ are the invariant
ratios. Cancellation occurs when a factor of $d_i$ divides $n_i$.
The topological invariant
(\Cref{def:winding}) generalises to the degree of the map
$z \mapsto \det H(z;\tau)$, or equivalently the sum of winding
numbers of the invariant ratios. The saddle-to-saddle structure
persists, but transitions can involve \emph{multiple simultaneous
activations} corresponding to rank jumps of $H$ by more than one.
\end{remark}

%======================================================================
\section{Summary of Predictions}\label{sec:predictions}
%======================================================================

\noindent
The framework generates the following precise, testable
predictions for recurrent network training under general data
distributions.

\begin{enumerate}[label=\textbf{P\arabic*.}]
  \item \textbf{Stepwise mode acquisition.} The effective rank
    of a recurrent network's transfer function increases
    as a step function during training, with each step
    corresponding to a single pole-zero separation event
    (\Cref{cor:staircase}). This holds for any smooth loss and
    any data distribution.

  \item \textbf{Data-dependent ordering.} The order of mode
    activation follows decreasing effective spectral energy
    $E_k^{\mathrm{eff}} = E_k^{\mathrm{target}} \cdot S_x(\arg p_k^*)$,
    the product of target spectral content and input spectral
    density at each mode's frequency (\Cref{cor:ordering}).
    Under white noise, this reduces to ordering by target
    energy alone (\Cref{cor:white}).

  \item \textbf{Input-dependent learning speed.} Modes at
    frequencies where the training input is energetic are
    activated faster; modes at frequencies absent from the input
    are activated slowly or not at all, regardless of target
    spectral content (\Cref{rem:structured}).

  \item \textbf{Dwell time scaling.}
    The plateau duration before the $(\rho+1)$-th mode activates
    scales as
    $T_\rho \sim \frac{1}{E_{\rho+1}^{\mathrm{eff}}}
    \log\frac{1}{|\eps_0|}$
    (\Cref{thm:dwell}).

  \item \textbf{Separation direction in $\C$.} At each
    transition, the pole-zero pair separates along a direction
    determined by the phase of the cross-correlation between the
    target residual and the input-filtered mode response. For
    structured inputs, this direction can differ from the phase
    of the target pole (\Cref{cor:geometry}).

  \item \textbf{Oscillatory transients.} When the driving
    coefficient $\lambda$ is complex, the transition is
    accompanied by a transient overshoot at the mode's frequency,
    with frequency and growth rate predicted by
    \Cref{thm:spiral}.

  \item \textbf{Topological robustness.} Early-activated modes
    (large pole-zero separation) are more robust to parameter
    perturbation than late-activated modes (small separation)
    (\Cref{cor:protection}). This is independent of the loss
    and data distribution.

  \item \textbf{Interference geometry.} Catastrophic forgetting
    under sequential training is predicted by the zero-trajectory
    geometry: a mode is forgotten if and only if a free zero,
    mobilised by the new task gradient, reaches the mode's pole
    in $\C$ (\Cref{thm:catastrophic}).

  \item \textbf{Data-overlap protection.} Modes whose frequencies
    carry effective spectral energy under both tasks are protected
    from catastrophic interference. Protection requires both
    input excitation and target demand at the mode's frequency
    (\Cref{cor:forgetting}).

  \item \textbf{Nonlinear dwell-time inflation.} In nonlinear
    recurrent networks, plateau durations are longer than the
    linear prediction by a factor controlled by the heterogeneity
    of instantaneous pole-zero separations across operating points
    (\Cref{prop:coherent}).
\end{enumerate}

%======================================================================
\appendix
\section{Proof Details for the Hessian Structure}\label{app:hessian}
%======================================================================

We provide the full computation of the Hessian at the cancellation
saddle in the general data-dependent setting.

Let $\theta = (W, b, c) \in \R^{n^2 + 2n}$ be the full parameter
vector. At a cancellation saddle where $r_k = \alpha_k\beta_k = 0$
with $\alpha_k = c^{\top}v_k = 0$ and $\beta_k = u_k^*b \neq 0$
(the case $\beta_k = 0$ is analogous by duality), we work in the
modal parameterisation $(\alpha_k, \beta_k)$ with all other
coordinates held fixed.

The loss restricted to $(\alpha_k, \beta_k)$ near the saddle,
for squared error under data distribution $\cD$, is
(from \eqref{eq:loss_expanded}):
\begin{align}
  \cL(\alpha_k, \beta_k)
  &= \cL_{\perp}
  + |\alpha_k|^2\,|\beta_k|^2\, S_k^{(x)}
  - 2\,\Re\!\bigl[\alpha_k\beta_k\, \bar{J}_k\bigr]
  + O(|\alpha_k\beta_k|^2),
\end{align}
where the mode energy $S_k^{(x)}$ and cross-energy $J_k$ are
defined in \Cref{def:mode_energy}, and
$\cL_{\perp}$ is the loss contribution from all other modes.

At $\alpha_k = 0$, $\beta_k = \beta_k^{(0)}$:
\begin{align}
  \frac{\partial\cL}{\partial\alpha_k} &= 0, \qquad
  \frac{\partial\cL}{\partial\beta_k} = 0, \\
  \frac{\partial^2\cL}{\partial\alpha_k\partial\bar\alpha_k}
  &= |\beta_k^{(0)}|^2\, S_k^{(x)} > 0, \\
  \frac{\partial^2\cL}{\partial\alpha_k\partial\bar\beta_k}
  &= -\bar{J}_k, \\
  \frac{\partial^2\cL}{\partial\beta_k\partial\bar\beta_k}
  &= 0.
\end{align}
The real Hessian restricted to
$(\Re\alpha_k, \Im\alpha_k, \Re\beta_k, \Im\beta_k)$
has eigenvalues (each with multiplicity 2):
\begin{equation}
  \mu_{\pm} = \frac{|\beta_k^{(0)}|^2 S_k^{(x)}}{2}
  \pm \sqrt{\frac{|\beta_k^{(0)}|^4 (S_k^{(x)})^2}{4}
  + |J_k|^2}.
\end{equation}
Since $S_k^{(x)} > 0$ (any nondegenerate input distribution
has positive energy in every mode) and $|J_k| > 0$
(hypothesis (c) of \Cref{thm:saddle}), we have $\mu_+ > 0$
and $\mu_- < 0$, confirming the saddle.

Note that the Hessian eigenvalues depend on the data
distribution through $S_k^{(x)}$ and $J_k$. The escape rate
$\mu_- \propto -E_k^{\mathrm{eff}}$ determines the
instability of the saddle: modes with higher effective spectral
energy produce more unstable saddles and faster escape.

%======================================================================
\section{Relationship to the Saxe--McClelland--Ganguli Framework}%
\label{app:saxe}
%======================================================================

We make explicit the formal correspondence between the feedforward
theory of Saxe et al.\ and the recurrent theory developed here.

\medskip
\begin{center}
\renewcommand{\arraystretch}{1.4}
\begin{tabular}{@{}p{0.42\textwidth}p{0.42\textwidth}@{}}
\hline
\textbf{Deep Linear Network (Saxe)} &
\textbf{Linear Recurrent Network (This work)} \\
\hline
End-to-end map $W_L\cdots W_1$ &
Transfer function $H(z;\tau) = C(zI-W)^{-1}B$ \\
Singular values $\sigma_k$ &
Residues $r_k$ at poles $p_k$ \\
Rank of $W_L\cdots W_1$ &
Effective rank $\rho = \#\{k:r_k\neq 0\}$ \\
Saddle: $\sigma_k = 0$ &
Saddle: $r_k = 0$ (pole-zero cancellation) \\
Activation: $\sigma_k: 0 \to {>}0$ &
Activation: pole-zero separation \\
Ordering by target SVD &
Ordering by $E_k^{\mathrm{eff}} = E_k^{\mathrm{target}}\cdot S_x$ \\
Dwell time $\sim 1/\sigma_k^{*2}$ &
Dwell time $\sim 1/E_k^{\mathrm{eff}}$ \\
Singular values $\in \R_{\ge 0}$ &
Residues $\in \C$ \\
No topological invariant &
Winding number $\cN(\gamma_k) \in \Z$ \\
No spiral transients &
Spiral separation for complex poles \\
No catastrophic interference mechanism &
Non-partner collision \\
Loss-independent ordering &
Input-dependent ordering \\
\hline
\end{tabular}
\end{center}

\medskip

\noindent
The central conceptual advance is the replacement of the singular
value (a non-negative real number) with the pole-zero configuration
(a pair of points in $\C$). The additional structure of
$\C$---the argument principle, residues, the geometry of
collisions in two dimensions---provides a richer and more
predictive theory. The data-dependent ordering through
$E_k^{\mathrm{eff}}$ is a genuinely new prediction: in
Saxe's framework, the singular values of the target determine
the ordering independently of the input statistics.

\end{document}
